\section{Simultaneous continuation of several orders}
\label{sec:simultaneous-covers}

Two individually finite covers need not have a common cover whose degree
is their least common multiple.  For the prime determinants of orders two
and three on \(\Hp_{1/5}\), the individual degrees are six and twelve.
The common continuation has degree twenty-four.  The additional sheets
record how the two branches change together, information that their
individual degrees do not retain.

\subsection{The vector of period characters}

Fix positive integers \(\mathbf m=(m_1,\ldots,m_r)\), a real number
\(\epsilon>0\), and \(D=\Hp_\epsilon\).  All functions and covers in
this section are complex analytic.  Their products are ordinary products
of scalar meromorphic functions, and their algebraic degrees refer to the
field \(\mathcal M(D)\).  No cohomological grading is assigned to them.
Put
\[
 J_\epsilon=\{n\in\mathbb Z_{\geq1}:n\epsilon<1\},\qquad
 v_n=(c_{m_1}(n),\ldots,c_{m_r}(n))\in\mathbb Q^r.
\]
Let \(S=\bigcup_i(S_{m_i}\cap D)\) and \(D^\circ=D\setminus S\).
Choose a real base point \(b>\max(1,\epsilon)\), with the positive
values \(\Delta_{m_i}(b)\) given by the initial products.  For a loop
\(\gamma\) in \(D^\circ\) based at \(b\), define its additive period
character by
\[
 \theta(\gamma)=
 \left(\frac1{2\pi i}\int_\gamma\omega_{m_i}\bmod\mathbb Z\right)_{i=1}^r
 \in(\mathbb Q/\mathbb Z)^r.
\]
Rationality follows from the residue formula.  Write
\(G_{\mathbf m}(\epsilon)=\operatorname{im}\theta\).
Its \(i\)-th coordinate, exponentiated by \(x\mapsto e^{2\pi ix}\),
is the character of the \(i\)-th determinant.

\begin{theorem}[Exact simultaneous character]
\label{thm:simultaneous-character}
One has
\[
 G_{\mathbf m}(\epsilon)
 =\left\langle v_n\bmod\mathbb Z^r:n\in J_\epsilon\right\rangle
 \subset(\mathbb Q/\mathbb Z)^r.
\]
This is a finite abelian group, which need not be cyclic.
The connected pointed cover corresponding to \(\ker\theta\) is the
minimal cover on which all the initial germs become single-valued.
It has degree \(|G_{\mathbf m}(\epsilon)|\), carries the functions
\(\mathcal D_{m_i}\), and admits a finite proper branched filling over
\(D\).  These are the least possible degree and the joint algebraic
degree of the germs over \(\mathcal M(D)\).

At \(a\in S\), form the total rational vector
\[
 \alpha(a)=(\alpha_{m_i}(a))_{i=1}^r
 =-\sum_{n\in J_\epsilon}\ord_a\zeta(ns)\,v_n.
\]
Let \(e(a)\) be the least positive integer for which
\(e(a)\alpha(a)\in\mathbb Z^r\).  The filling has
\(|G_{\mathbf m}(\epsilon)|/e(a)\) points above \(a\).
Each has ramification index \(e(a)\), and the \(i\)-th determinant
has order \(e(a)\alpha_{m_i}(a)\) there.
\end{theorem}

\begin{proof}
Choose \(N\geq\max_i m_i\) with \(1/(N+1)\leq\epsilon\).
In the extraction formula, every factor with \(n\epsilon\geq1\)
has its original holomorphic logarithm \(Z(ns)\) on \(D\).
Absorb these terms into the holomorphic functions
\[
 h_i(s)=T_{m_i,N}(s)
       -\sum_{\substack{1\leq n\leq N\\n\epsilon\geq1}}
          c_{m_i}(n)Z(ns).
\]
Terms with zero coefficient are omitted.  Then
\[
 \omega_{m_i}=dh_i-
       \sum_{n\in J_\epsilon}c_{m_i}(n)\,d\log\zeta(ns).
\]
The logarithmic differential is single-valued even where the logarithm
is not.  Summing its orders at each point gives the displayed vector
\(\alpha(a)\), including every coincident zero.  The residue-period
formula on the convex half-plane \(D\) shows that every period is an
integer combination of these vectors.  Thus the image is contained in
the group generated by the \(v_n\).

For the reverse inclusion, \(1/n\in D\) exactly when
\(n\epsilon<1\).  At \(1/n\), only \(\zeta(ns)\) has a pole.
Every other scaled argument is positive real and different from one,
so its zeta value is finite and nonzero.  A small positive loop about
\(1/n\), connected to \(b\) by any path in \(D^\circ\), therefore
has period vector \(v_n\bmod\mathbb Z^r\).  Zero vectors need no
loop.  This proves equality without any assumption on complex zeta zeros.
There are finitely many rational generators, so the group is finite.

Use paths from \(b\), identifying two with the same endpoint exactly
when the integrals of every \(\omega_{m_i}\) differ by
\(2\pi i\mathbb Z\).  The disk charts of
Theorem~\ref{thm:minimal-cover} give a connected cover, and on it put
\[
 \mathcal D_{m_i}[\gamma]
 =\Delta_{m_i}(b)\exp\int_\gamma\omega_{m_i}.
\]
The fibre is a torsor for the vector image.  A pointed cover carrying
all these functions maps to this one by projected paths; the functions
make the map well-defined, and disk charts prove holomorphy and uniqueness.
Distinct period vectors give distinct tuples of nonzero values above
\(b\).  Any finite realization of all germs must therefore have at
least \(|G_{\mathbf m}(\epsilon)|\) sheets.

Near \(a\), write
\(\omega_{m_i}=\alpha_{m_i}(a)\,ds/(s-a)+dq_i\), with \(q_i\)
holomorphic on a disk.  The local sheets form cycles of length \(e(a)\).
Fill each cycle by
\[
 s=a+t^{e(a)},\qquad
 \mathcal D_{m_i}=C_i t^{e(a)\alpha_{m_i}(a)}
                    e^{q_i(a+t^{e(a)})},\qquad C_i\ne0.
\]
The constants are fixed by the continued germs.  These charts give all
points and orders stated in the theorem.  Changing a root coordinate
multiplies \(t\) by a root of unity, so the charts glue to the
punctured cover.  They also prove properness: the inverse image of a small
closed disk is a finite union of compact disks, and finitely many such
base disks cover any compact subset of \(D\).
Any meromorphic local realization requires all its pulled-back orders
to be integers, hence a local index divisible by \(e(a)\).

The functions on this finite cover are algebraic over \(\mathcal M(D)\):
the product of \(X-F\) over its sheets has single-valued meromorphic
coefficients for each meromorphic \(F\).  Every finitely generated subextension of its meromorphic function
field is finite algebraic.  In characteristic zero it has a primitive
element: for two generators, all but finitely many complex linear
combinations separate their finitely many embeddings, and induction
handles any finite list.  The sheet polynomial bounds the degree of
that primitive element by the cover degree.  Thus the whole function
field has degree at most the cover degree: any larger linearly
independent finite list would lie in such a subextension.
For a direct verification of equality for the field generated by our
coordinates, choose complex constants \(\lambda_i\) so that the values
of \(F=\sum_i\lambda_i\mathcal D_{m_i}\) are distinct on the fibre
over \(b\).  This excludes finitely many proper linear hyperplanes,
since the tuples on that fibre are distinct.  The continued germ of
\(F\) has exactly \(|G_{\mathbf m}(\epsilon)|\) distinct values there.
Every polynomial annihilating it has at least that degree, as in the
proof of Theorem~\ref{mono:thm:criterion}.  Its sheet product attains
that degree.  The coordinate field contains \(F\) and lies in the
function field of the cover, proving the asserted joint degree.
\end{proof}

The strict cutoff in this proof is essential.  If \(n\epsilon=1\),
then \(\operatorname{Re}(ns)>1\) everywhere in \(D\); its pole is on
the excluded boundary, and its logarithm contributes no period in \(D\).
At a common complex zero, the sum defining \(\alpha(a)\) must be taken
before denominators are reduced.  Thus coincident divisors can reduce
local ramification, while the separate reciprocal-real poles already
force the full group in the theorem.

\subsection{Integer lattices and coordinates}

The theorem turns the degree into a finite integer calculation.
Choose a positive integer \(Q\) divisible by the denominators of every
coordinate of the finitely many \(v_n\).  Let \(A\) be the integer
matrix with columns \(Qv_n\), and define the full-rank lattice
\[
 \Lambda=Q\mathbb Z^r+A\mathbb Z^{J_\epsilon}\subset\mathbb Z^r.
\]
Here \(\mathbb Z^{J_\epsilon}\) is the free group with one generator
per index in \(J_\epsilon\).  Reduction modulo \(Q\mathbb Z^r\)
gives
\[
 G_{\mathbf m}(\epsilon)\cong\Lambda/Q\mathbb Z^r,
 \qquad
 |G_{\mathbf m}(\epsilon)|=\frac{Q^r}{[\mathbb Z^r:\Lambda]}.
\]
One exact procedure is integer column reduction of \([QI_r\ A]\)
to a basis matrix \(B\) for \(\Lambda\).  Then
\(C=B^{-1}QI_r\) is integral, and
\[
 G_{\mathbf m}(\epsilon)\cong\mathbb Z^r/C\mathbb Z^r.
\]
Apply integer row and column operations of determinant \(\pm1\) to
\(C\), obtaining its Smith diagonal
\(d_1,\ldots,d_r\), with \(d_i>0\) and \(d_i\mid d_{i+1}\).
The group is \(\bigoplus_i\mathbb Z/d_i\mathbb Z\), with factors
of order one omitted, and its degree is \(\prod_i d_i\).

These reductions use only the Euclidean algorithm.  To recall why they
terminate, move a nonzero entry to the first position and use division
with remainder to reduce its row and column until it divides them.
If it fails to divide an entry in the remaining block, adding that
entry's row to the first row and repeating produces a smaller positive
pivot.  Eventually the pivot divides every remaining entry.  Clear its
row and column and repeat in the smaller block.  The resulting diagonal
has the stated divisibility.  Every operation changes a basis in the
source or target free group, so it preserves the quotient.  Applied to
the rectangular generating matrix, the same operations produce a lattice
basis.  This proves the procedure and its degree formula, including
redundant or zero generators.

For explicit meromorphic coordinates put
\[
 k_i=M_{m_i}(\epsilon),\qquad
 A_i(s)=e^{k_i h_i(s)}
        \prod_{n\in J_\epsilon}\zeta(ns)^{-k_i c_{m_i}(n)}.
\]
The powers are integers.  Over \(D^\circ\), form the product of the
root equations
\[
 y_i^{k_i}=A_i(s)\quad(1\leq i\leq r),\qquad
 y_i(b)=\Delta_{m_i}(b).
\]
Coordinatewise multiplication by roots of unity gives an action of
\(G_0=\prod_i(\tfrac1{k_i}\mathbb Z/\mathbb Z)\) on the product,
simply transitive on every fibre. Continuation acts by translations
in the subgroup \(G_{\mathbf m}(\epsilon)\).
Consequently the product has
\(\prod_i k_i/|G_{\mathbf m}(\epsilon)|\) connected components.
The component containing the displayed base point is the minimal
simultaneous cover.  Closing the equations in
\(D\times(\mathbb P^1)^r\) and replacing each local branch by the
smooth coordinates of Theorem~\ref{thm:simultaneous-character} gives
its normalization and branched filling.

The selected component also has finite equations beyond the individual
root equations.  A character of \(G_0\) is represented by integers
\(a_1,\ldots,a_r\), with \(a_i\) taken modulo \(k_i\), and has
value \(\sum_i a_i g_i\bmod\mathbb Z\) at \(g\in G_0\).
It vanishes on the simultaneous subgroup exactly when
\[
 d_n(a):=\sum_i a_i c_{m_i}(n)\in\mathbb Z
 \quad\text{for every }n\in J_\epsilon.
\]
For each such character define
\[
 R_a(s)=\exp\left(\sum_i a_i h_i(s)\right)
          \prod_{n\in J_\epsilon}\zeta(ns)^{-d_n(a)}.
\]
It is meromorphic on \(D\).  The equations
\(\prod_i y_i^{a_i}=R_a(s)\) select the pointed component.
Indeed, on any disk the root tuples differ by \(G_0\), and these
equations retain exactly the common kernel of the characters vanishing
on \(G_{\mathbf m}(\epsilon)\).  That kernel is the subgroup itself:
Smith reduction of the quotient \(G_0/G_{\mathbf m}(\epsilon)\)
provides a character detecting every nonzero quotient element.
A finite generating set of its character group suffices.
Changing \(a_i\) by \(k_i\) multiplies both sides by \(A_i\), so
this description is independent of the chosen representatives.

Coordinate projection maps the minimal simultaneous cover onto the
minimal \(m_i\)-cover, with degree \(|G_{\mathbf m}(\epsilon)|/k_i\).
More generally, deleting coordinates gives the cover for their projected
period group.  In the local coordinates at a divisor, such a map is
\(t\mapsto t^{e/e'}\), up to the root of unity fixed by the base
point, where \(e'\) is the local period order after projection.
This also proves extension over every filled point.

The order-change multiplier is the meromorphic function \(y_j/y_i\).
Its period character is \(g\mapsto g_j-g_i\), its logarithmic differential is
\(\omega_{m_j}-\omega_{m_i}\), and its local order on the simultaneous
cover is \(e(a)(\alpha_{m_j}(a)-\alpha_{m_i}(a))\).
The ratios satisfy the cocycle identity by cancellation.
Its minimal cover is the quotient by the kernel of this difference
character.  A map from the individual \(i\)-cover to the individual
\(j\)-cover preserving continued paths exists exactly when
\(\ker(g\mapsto g_i)\subseteq\ker(g\mapsto g_j)\).
Thus the scalar descent condition is retained, while the common cover
is governed by the whole vector group.

If \(\epsilon'>\epsilon\), then \(J_{\epsilon'}\subseteq J_\epsilon\)
and the smaller half-plane has the subgroup generated by the smaller
list.  Restricting the larger continuation surface gives components
indexed by \(G_{\mathbf m}(\epsilon)/G_{\mathbf m}(\epsilon')\).
The positive initial branches select one component and its compatible
inclusion.  For a connected domain with holes, the same statement uses
its actual period subgroup; the finite list of divisor points lying
inside the domain alone does not determine it.

\subsection{Orders two and three on \texorpdfstring{\(\Hp_{1/5}\)}{Re(s) > 1/5}}
\label{subsec:simultaneous-two-three}

Here \(J_{1/5}=\{1,2,3,4\}\).  The nonzero period generators are
\[
 g_2=(\tfrac12,0),\qquad
 g_3=(\tfrac13,\tfrac13),\qquad
 g_4=(0,\tfrac14).
\]
They have orders two, three and four.  If
\(a g_2+b g_3+c g_4=0\), with \(0\leq a<2\), \(0\leq b<3\),
\(0\leq c<4\), the first coordinate gives
\(3a+2b\equiv0\pmod6\).  Reduction modulo two and three forces
\(a=b=0\); the second coordinate then forces \(c=0\).
Therefore
\[
 G_{(2,3)}(1/5)
 \cong\mathbb Z/2\mathbb Z\oplus\mathbb Z/3\mathbb Z
                    \oplus\mathbb Z/4\mathbb Z
 \cong\mathbb Z/2\mathbb Z\oplus\mathbb Z/12\mathbb Z.
\]
The first coordinate image has order six and the second has order
twelve.  Their least common multiple, twelve, is strictly smaller than
the simultaneous degree twenty-four.

The lattice calculation gives the same invariants by a different route.
With \(Q=12\), the generator columns are \((6,0),(4,4),(0,3)\).
Together with \(12\mathbb Z^2\), they span
\[
 \Lambda=\{(x,y)\in\mathbb Z^2:x\equiv0\pmod2,
                                      \ x-y\equiv0\pmod3\}.
\]
A basis and its relation matrix are
\[
 B=\begin{pmatrix}6&4\\0&1\end{pmatrix},\qquad
 B^{-1}(12I_2)=\begin{pmatrix}2&-8\\0&12\end{pmatrix}.
\]
The greatest common divisor of the latter entries is two and its
determinant is twenty-four.  Its Smith invariants are consequently two
and twelve.  The subgroup consists of the twenty-four pairs
\((a/6,b/12)\) with \(b-2a\equiv0\pmod3\).

To see the functions, abbreviate \(T_2=T_{2,4}\), \(T_3=T_{3,4}\),
and \(z_n(s)=\zeta(ns)\).  On the simultaneous cover there are roots
\[
 u^2=z_2^{-1},\qquad v^3=z_3^{-1},\qquad w^4=z_4^{-1},
\]
chosen positive at \(s=2\).  Their product cover is connected: the
loops about \(1/2,1/3,1/4\) change precisely one of these three
roots and generate all \(2\cdot3\cdot4\) sheets.  The determinant
coordinates are
\[
 y_2=e^{T_2}uv,\qquad y_3=e^{T_3}vw.
\]
No two root triples have the same pair.  For a ratio of two triples,
\(uv=vw=1\) and \(u^2=v^3=w^4=1\) force \(v=u=w=1\).
One can also recover the functions directly: writing
\(q_2=e^{-T_2}y_2\), \(q_3=e^{-T_3}y_3\), one has
\[
 v=\frac1{z_2z_3q_2^2},\qquad u=q_2/v,\qquad w=q_3/v.
\]
These identities are meromorphic and continue across the filled points.
Thus the root triples and the determinant pair describe the same
normalized connected surface, including its functions.

The individual equations are
\[
 y_2^6=A:=e^{6T_2}z_2^{-3}z_3^{-2},\qquad
 y_3^{12}=B_3:=e^{12T_3}z_3^{-4}z_4^{-3}.
\]
Their fibre product has seventy-two sheets before a component is
selected.  Put
\[
 C=e^{4(T_3-T_2)}z_2^2z_4^{-1}.
\]
Then \(C^3=B_3/A^2\), and the three components are distinguished by
\[
 y_3^4=\xi C y_2^4,\qquad \xi^3=1.
\]
The positive values at \(2\) select \(\xi=1\).  On that component
\(y_2\) and \(y_3\) have the specified germs; its projections have
degrees four and two.  The normalization is therefore essential to the
geometric comparison, beyond the separate degrees of the equations.

The local data at the reciprocal-real poles are
\[
\begin{array}{c|c|c|c|c}
 a&\alpha(a)&e(a)&\text{points above }a&
        (\ord y_2,\ord y_3)\\\hline
 1/2&(1/2,0)&2&12&(1,0)\\
 1/3&(1/3,1/3)&3&8&(1,1)\\
 1/4&(0,1/4)&4&6&(0,1)
\end{array}
\]
At any coincident complex zeros, the summed vector in
Theorem~\ref{thm:simultaneous-character} gives the corresponding
ramification and orders without a simplicity assumption.

Finally \(H_{2,3}=y_3/y_2\) has character values generated by
\(-1/2\) and \(1/4\), hence a four-sheeted minimal cover and a kernel
of order six on the simultaneous cover.  Its fourth power is \(C\).
The loop at \(1/4\) is trivial on \(y_2\) and nontrivial on \(y_3\),
whereas the loop at \(1/2\) is trivial on \(y_3\) and nontrivial on
\(y_2\).  Neither individual minimal cover maps to the other while
preserving continued paths.  The simultaneous cover supplies both
projections and their multiplier with the least possible degree.
This comparison concerns scalar determinant germs.  A map to an
arithmetic complex would additionally have to preserve its differential,
products and arithmetic classes; none is supplied by these root equations.
